\documentclass[12pt]{article}
\usepackage{latexblog}

% ============================================================
% Blog Metadata
% ============================================================
\blogtitle{Scheme语言的hello world}
\blogdate{2021-01-05}
\blogtags{编程语言, Scheme, 闲话编程}
\bloglang{zh}
\blogsummary{}

\begin{document}
\makeblogtitle

最近准备创建一个简单的基于JVM的新语言，但在语法设计上犹豫不决，研究了一些资料后，觉得可能Lisp正是我所追求的清晰、简单、优雅的语法参考，因此决定深入了解一下Scheme。

\section{开发环境配置}\label{ux5f00ux53d1ux73afux5883ux914dux7f6e}

\subsection{运行环境}\label{ux8fd0ux884cux73afux5883}

编译安装\href{https://github.com/cisco/ChezScheme}{ChezScheme}，据说是最好的scheme实现。

\begin{Shaded}
\begin{Highlighting}[]
\FunctionTok{git}\NormalTok{ clone https://github.com/cisco/ChezScheme.git}
\BuiltInTok{cd}\NormalTok{ ChezScheme}
\ExtensionTok{./configure}
\FunctionTok{make}
\FunctionTok{sudo}\NormalTok{ make install}
\end{Highlighting}
\end{Shaded}

完成之后应该就可以来写一个hello world了：

\begin{verbatim}
Chez Scheme Version 9.5.5
Copyright 1984-2020 Cisco Systems, Inc.

> (+ 1 1)
2
>
> (display "Hello world")
Hello world
\end{verbatim}

\subsection{Visual studio code}\label{visual-studio-code}

需要安装两个插件

\begin{itemize}
\tightlist
\item
  vscode-scheme
\item
  Code Runner
\end{itemize}

然后在settings中查找''Code-runner: Executor Map''，修改其中的scheme命令行：

\begin{Shaded}
\begin{Highlighting}[]
\StringTok{"code{-}runner.executorMap"}\ErrorTok{:} \FunctionTok{\{}
    \DataTypeTok{"scheme"}\FunctionTok{:} \StringTok{"scheme {-}{-}script"}
\FunctionTok{\}}
\end{Highlighting}
\end{Shaded}

参考：

\begin{itemize}
\tightlist
\item
  \href{https://www.yinwang.org/blog-cn/2013/04/11/scheme-setup}{Scheme 编程环境的设置}
\end{itemize}


\end{document}
